\chapter{Literature Review}

This chapter covers the research that shaped our project. The goal is not to list tools or summarize papers, but to explain what the existing literature taught us and where it left questions that our project tries to answer.

The research was divided among the team. Arihant and Huizhe took on distributed systems theory: consistency models, replication strategies, and synchronization algorithms. They focused on the academic side, reading papers by Lamport, Shapiro, and Saito. Nizar and Shengkang surveyed existing tools and local networking technologies, testing products like Syncthing and comparing cloud platforms. Each week, we met to present what we found and argue about what mattered for our design. Some of the best design decisions came out of those arguments.

\section{The Collaborative File Management Landscape}

Collaborative file management tools can be grouped into four categories. None of them does everything well, and understanding where each one falls short is what pointed us toward our project.

\textbf{Centralized cloud platforms} (Google Drive \cite{googledrive}, Dropbox \cite{dropbox}, Nextcloud) are the tools most students reach for first. They synchronize files through a central server, handle multiple devices, and have friendly interfaces. The problems are structural, not cosmetic. Everything depends on the server. If the internet is down, or if the campus Wi-Fi drops, you cannot sync. If the provider changes their pricing or storage limits, you adjust or leave. Dropbox's own engineering team rebuilt their sync engine from scratch to handle the complexity of tracking file states across devices \cite{dropboxtech}, which illustrates how difficult reliable synchronization is even with centralised infrastructure. Version history exists, but it is time-limited (30 days on Google Drive, for example) and not designed for tracking the evolution of a project.

\textbf{Version control systems} (Git \cite{gitbook}, SVN, Perforce) solve the history problem completely. Git stores every version of every file, supports branching, and has detailed merge conflict resolution. But Git was built for software developers. The staging-committing-pushing workflow is not intuitive for someone who just wants to share a presentation or a dataset. In our own team, two out of four members had never used Git before this project, which is exactly the audience we are designing for.

\textbf{Peer-to-peer synchronization tools} (Syncthing \cite{syncthing}, Resilio Sync \cite{resilio}) take a different approach. They synchronize files directly between devices, no server needed. Syncthing in particular is open source, works on most platforms, and handles any file type. But when Nizar tested it for a week, he noticed a gap: if he accidentally saved a file with errors, it immediately replaced the file on all connected devices. Syncthing does have a versioning feature, but it is not on by default and the interface for recovering old versions is buried in the settings. Versioning is an afterthought, not part of the core design.

\textbf{Local sharing} (AirDrop, Nearby Share) is the simplest category. Select a file, tap a device, done. But these are one-shot transfers. There is no persistent connection, no synchronization, and no history. They solve a different problem.

The pattern is clear: each category does well on some dimensions and poorly on others. No single tool provides automatic synchronization, version safety, and zero configuration at the same time. The detailed comparison tables in Appendix~\ref{app:comparison-tables} show this across five evaluation dimensions.

\section{Synchronization in Distributed Systems}

Rather than providing a comprehensive survey of distributed systems literature, this section highlights the specific synchronization challenges and design trade-offs that guided our implementation.

\subsection{Ordering Events Without a Shared Clock}

The first problem any distributed system faces is ordering. When things happen on different machines, how do you decide what came first? Physical clocks are unreliable because machines drift.

Lamport \citeyear{lamport1978} proposed a simple solution: logical clocks. Each process keeps a counter and increments it on every event. When sending a message, the sender attaches its counter. The receiver updates its own counter to be at least as high. This creates a \emph{partial ordering}: if event $A$ causally led to event $B$, then $A$'s timestamp will always be lower.

The word ``partial'' matters. Lamport clocks cannot tell you whether two events are concurrent, that is, whether they happened independently with no causal connection. Fidge \citeyear{fidge1988} and Mattern independently proposed vector clocks to detect concurrency. Each node keeps an array of counters, one per participant. Vector clocks can identify concurrent events, but the counter array grows with the number of peers. Earlier, Parker et al. \citeyear{parker1983} had introduced version vectors for detecting mutual inconsistency between replicas, the formal basis for recognising when two copies of the same data have diverged.

For our system, this trade-off was the first real design decision. With 3 to 5 peers on a local network, vector clocks would mean carrying a small array of integers. The overhead is tiny. But vector clocks also require knowing all participants, and adding or removing a peer requires updating every node's vector. Since our system is designed for temporary groups where members join and leave, Lamport clocks are simpler to manage. The cost is that we cannot detect true concurrency. If two peers edit the same file at the exact same moment, one version will silently overwrite the other. The backup mechanism makes this acceptable, but it is a real limitation.

\subsection{The CAP Theorem and Why It Matters}

Gilbert and Lynch \citeyear{brewer2002} proved that a distributed system cannot simultaneously guarantee consistency, availability, and partition tolerance. During a network partition (which on a campus Wi-Fi happens every time a laptop lid closes), you must choose: either block operations until the partition heals (consistency), or allow operations to continue and reconcile later (availability).

In our case, we have to prioritize availability. Student laptops disconnect constantly: Wi-Fi drops, people move between classrooms, or a battery dies. If the sync tool blocked edits whenever a teammate was offline, it would be unusable. We therefore accept eventual consistency, allowing each peer to edit files offline and sync changes once they reconnect.

Vogels \citeyear{vogels2009} argued that eventual consistency is the right model for applications where reading slightly stale data is harmless. A teammate seeing a file that is ten seconds out of date causes no damage. Losing a file because the system refused to save it during a disconnection does.

\subsection{Resolving Conflicts}

When two peers edit the same file while disconnected, the system must decide what happens when they reconnect. The literature offers several approaches, each with trade-offs.

\textbf{Operational transformation} (OT) \cite{sun1998} tracks individual editing operations (insert character at position 5, delete characters 3 through 7) and transforms them so they can be applied in any order. Google Docs relies on this. OT is powerful for text, but it requires operation-level logging. For binary files like images, spreadsheets, or compiled code, there are no ``operations'' to transform.

\textbf{CRDTs} (Conflict-free Replicated Data Types) \cite{shapiro2011} are data structures that mathematically guarantee convergence. Shapiro et al. proved that any data type whose merge operation is commutative, associative, and idempotent will converge regardless of operation order. This is a strong result. But applying CRDTs to file metadata requires defining what ``merge'' means for each field in the metadata schema, and testing the merge logic thoroughly.

\textbf{Last-Write-Wins} (LWW) is the simplest policy: keep the version with the highest timestamp. One version is discarded. The Bayou project \cite{bayou1995} demonstrated that LWW becomes safe when combined with backup storage: before overwriting a file, save the old copy somewhere. Terry et al. \citeyear{terry1994} added ``session guarantees,'' ensuring that a user always sees the results of their own writes, even if other peers' updates have not arrived.

After reading these papers, the team debated for two meetings before settling on LWW with backups. CRDTs would be technically better, but implementing and testing them correctly within our timeline was unrealistic. Saito and Shapiro's \citeyear{saito2005} survey of optimistic replication confirmed that our approach, optimistic local updates with deferred reconciliation and rollback, is a well-documented pattern for weakly connected systems. That reassured us. The metadata schema we designed can support CRDTs in a future version, using register-based merge semantics from Shapiro et al. \citeyear{shapiro2011}, so the option remains open.

\subsection{Propagating Updates Efficiently}

Tridgell and Mackerras \citeyear{rsync1996} showed that rolling checksums allow two machines to identify which \emph{blocks} of a file have changed and transfer only the differences. This is the core idea behind rsync, and it reduces bandwidth use dramatically for large files with small changes.

Demers et al. \citeyear{demers1987} studied a separate problem: how to reliably spread updates across a group of peers without a central coordinator. Their ``epidemic'' (gossip) algorithms model updates as a contagion that spreads from peer to peer until everyone has it. The approach is probabilistic but robust, and it scales well.

For our initial version, we transfer whole files. It is simpler to implement, simpler to debug, and for files under a few megabytes on a local network, the performance penalty is small. The architecture is designed to support both optimizations later: block-level transfers in the C++ daemon, and gossip propagation in the Go coordinator.

\section{Identified Gap}

Comparing existing tools across five dimensions (membership management, synchronization capability, history management, file type support, and filesystem transparency) shows a consistent pattern. The full tables are in Appendix~\ref{app:comparison-tables}; here is the summary:

\begin{itemize}
    \item Cloud platforms score well on synchronization and usability, but poorly on history management and offline use.
    \item Version control systems score well on history, but poorly on synchronization automation and accessibility for non-technical users.
    \item P2P sync tools score well on synchronization and usability, but poorly on history management.
    \item Local sharing tools score well on usability but offer no synchronization or history at all.
\end{itemize}

The gap is not that any single feature is missing from the market. It is that no tool scores well across all five dimensions simultaneously. Each existing tool was optimised for a different audience (enterprise workers, developers, power users, casual sharers), leaving a gap for small, temporary teams who want something that ``just works'' on a local network with version safety included.

\section{Positioning the Project}

Our system targets this gap. It differs from existing tools in four concrete ways:

\begin{itemize}
    \item \textbf{LAN-only, serverless operation.} Peer discovery uses mDNS/DNS-SD on the local subnet. No servers, accounts, or internet connection needed.
    \item \textbf{Zero configuration.} Peers appear automatically. There is nothing to install per-project, no folder IDs to share, no repositories to initialise.
    \item \textbf{Automatic version backups.} Every overwrite saves the old file first. This is built into the synchronization protocol, not a feature that must be turned on.
    \item \textbf{Low-resource split architecture.} Go handles networking (peer coordination, TCP connections). C++ handles file I/O (watching, hashing, disk writes). This keeps CPU and memory use low on student laptops.
\end{itemize}

Three design principles guide the project:

\begin{enumerate}
    \item \textbf{Transparency.} The synchronized folder should look and feel like a normal local folder. No special commands. No merge dialogs.
    \item \textbf{Safety.} If the system changes a file, the old version must be recoverable. No automation should cause permanent data loss.
    \item \textbf{Lightness.} The tool runs continuously in the background. It must not drain battery life or make the laptop feel slow.
\end{enumerate}

\section{Summary}

The literature review taught us three things that directly determined the design of this project.

The CAP theorem says we must accept eventual consistency. Students disconnect constantly, and a synchronization tool that stops working when someone closes their laptop is useless. The system has to let everyone work offline and sort out differences later.

For conflict resolution, the right choice depends on scale and timeline. LWW with backups gives us simple, predictable behaviour with no risk of permanent data loss. CRDTs would be better for a production system, but they are complex to implement correctly. The schema supports adding them later.

No existing tool occupies the design space we are targeting. Cloud services, version control, P2P sync, and local sharing each solve part of the problem. None of them solve all of it at once, for a small team, on a local network, with no setup. That is where our project fits.
